% Revised November 20, 2026.
\documentclass[11pt]{amsart}

\usepackage[margin=1.15in]{geometry}
\usepackage[T1]{fontenc}
\usepackage{lmodern,amssymb,mathtools,microtype}
\usepackage{etoolbox,xcolor}
\usepackage[colorlinks=true,linkcolor=blue!55!black,citecolor=green!35!black,
  urlcolor=blue!55!black,pagebackref=true]{hyperref}
\hypersetup{pdftitle={Regularization of plurisubharmonic exhaustions on complex surfaces},
  pdfauthor={Zhijian Chen and Bo Xiao},
  pdfsubject={A proposed proof for smooth and normal complex surfaces}}
\renewcommand*{\backref}[1]{}
\renewcommand*{\backrefalt}[4]{%
  \ifnum#1>0\relax\ensuremath{\uparrow}\,#2\fi}
\renewcommand*{\backrefsep}{, }
\renewcommand*{\backreftwosep}{, }
\renewcommand*{\backreflastsep}{, }

\newtheorem{theorem}{Theorem}[section]
\newtheorem{lemma}[theorem]{Lemma}
\newtheorem{proposition}[theorem]{Proposition}
\newtheorem{corollary}[theorem]{Corollary}
\numberwithin{equation}{section}
\newcommand{\D}{\mathrm{D}}
\newcommand{\ddc}{dd^c}
\newcommand{\R}{\mathbb R}
\newcommand{\C}{\mathbb C}
\DeclareMathOperator{\supp}{supp}
\DeclareMathOperator{\osc}{osc}
\title[Regularization of plurisubharmonic exhaustions]{Regularization of
plurisubharmonic exhaustions on complex surfaces}
\author{Zhijian Chen}
\address{Institute of Mathematics, Academy of Mathematics and Systems Science, Chinese Academy of Sciences, Beijing 100190, China}
\email{chenzhijian@amss.ac.cn}
\author{Bo Xiao}
\address{School of Mathematical Sciences, University of Chinese Academy of Sciences, Beijing, China}
\email{xb1994@amss.ac.cn}
\date{November 20, 2026}
\subjclass[2020]{Primary 32U10; Secondary 32T35, 32U25, 32U40}
\keywords{Plurisubharmonic exhaustion, weak completeness, complex surface,
Hermitian metric, Lelong number, Levi-flat hypersurface}

\setcounter{tocdepth}{1}
\makeatletter
\patchcmd{\@starttoc}{\addvspace{32\p@\@plus14\p@}}{}{}%
  {\PackageError{surface-regularization}{Cannot adjust contents spacing}{}}
\makeatother

\begin{document}
\begin{abstract}
We give an intrinsic argument for smooth regularization on a complex
surface carrying a continuous plurisubharmonic exhaustion. Outside the
case of a modification of a Stein space, the intrinsic version of
the level average from our companion paper produces a nonzero compactly supported current of
uniformly vanishing two-dimensional density. A Hermitian heat operator
and a scalar Dirichlet problem realize this current as the complex Hessian
of a negative potential on arbitrarily large sublevels, with boundary value
zero. Regular levels of those potentials give Levi-flat boundaries. A
second Dirichlet problem and a surface intersection identity give a
smooth global exhaustion. Choosing the boundary levels away from
exceptional fibers extends the conclusion to normal complex surfaces.
\end{abstract}
\maketitle
\tableofcontents

\section{Introduction and main results}

In 1980 S.~Nakano asked whether every complex manifold admitting a
continuous plurisubharmonic (psh) exhaustion admits a smooth one
\cite[\S I.1, Problem~1, p.~162]{Nakano}. In the companion paper
\cite[Theorems~1.1 and 1.3]{CX} we construct a smooth quasi-projective
threefold with a locally Lipschitz psh exhaustion and no $C^1$ psh
exhaustion, and prove a positive result for divisor complements in normal
compact K\"ahler surfaces. We consider the surface question without a
compactification or a K\"ahler metric.

\begin{theorem}\label{thm:main}
Let \(X\) be a connected smooth complex manifold of complex dimension two.
If \(X\) admits a continuous real-valued plurisubharmonic exhaustion,
then \(X\) admits a \(C^\infty\) plurisubharmonic exhaustion.
\end{theorem}

We use exhaustion in the sense that all finite sublevels are relatively
compact. If the continuous input takes the extended value \(-\infty\),
replace it by \(\max(\rho,0)\) to obtain the real-valued hypothesis.
We write \(d^c=\tfrac i2(\bar\partial-\partial)\),
\(\D=dd^c=i\partial\bar\partial\). A Hermitian form \(\omega\) is fixed
on \(X\), \(dV=\omega^2/2\), and \(L=\Lambda_\omega\D\) has the sign of
the usual sum of second derivatives. The metric may have torsion.
In real dimension four, a real \((1,1)\)-form \(B\) with
\(\Lambda_\omega B=0\) satisfies \(B\wedge B=-|B|_\omega^2dV\).
This identity enters twice, first in recovering the current from a scalar
potential, and then in showing that harmonic functions between Levi-flat
boundaries are pluriharmonic.

The output is a new exhaustion. Since every connected noncompact Riemann
surface is Stein, Theorem~\ref{thm:main} and the threefold in \cite{CX}
put the smallest dimension of a smooth counterexample at three.
Corollary~\ref{cor:normal} treats normal surface singularities.
The nonnormal example in \cite[Proposition~5.7]{CX} shows
where that extension stops.

Several positive cases were known. For an irreducible noncompact
surface with a continuous psh exhaustion and a nonconstant global
holomorphic function, Tomassini and V\^aj\^aitu \cite{TV} prove
holomorphic convexity, giving a smooth psh exhaustion by pullback
from the Stein reduction. This extends Ohsawa's smooth case
\cite[Proposition~1.4]{Ohsawa}. V\^{\i}j\^{\i}itu
\cite[Theorem~1.1]{Vij} allows an arbitrary psh exhaustion in this
criterion. Mongodi and S{\l}odkowski \cite[Corollary~4.5]{MS} obtain
a real-analytic psh exhaustion for a continuously exhausted domain
in a surface already carrying a real-analytic psh exhaustion.

The proof starts with the level averaging of our companion paper
\cite[Lemma~5.3]{CX} in the intrinsic form of the October~29 project
note \cite[Proposition~2.1]{Oct}. That note also develops a heat
estimate on complete K\"ahler surfaces of bounded curvature and a
Green-potential construction under additional hypotheses at infinity.
Here a local Hermitian heat operator supplies an almost positive
representative with small $dd^c$ error (Lemma~\ref{lem:heat}).
On smooth outer approximations to a continuous sublevel, we control
the boundary term in the trace-correction argument of Buchdahl
\cite[\S2]{Buchdahl}. This produces a potential with prescribed
boundary limit (Lemma~\ref{lem:dirichlet}). We then apply the
boundary identity used by Weber \cite{Weber} to the Hermitian trace
equation on overlapping annuli and glue their pluriharmonic solutions.

\subsection*{Acknowledgement and AI disclosure}
OpenAI Codex developed and drafted the proposed Hermitian and Dirichlet
arguments in this paper and carried out the subsequent internal proof and
attribution review. The October~29 research note is cited under its
recorded title and byline. Independent mathematical review of this
preprint is pending.

\section{A current on a compact band of levels}

\begin{lemma}\label{lem:shell}
Suppose that \(X\) is noncompact, that \(\rho\) is a continuous psh
exhaustion, and that \(X\) is not a modification of a Stein space. There
exists a nonzero positive \((1,1)\)-current \(T\) of compact support
with \(\D T=0\). Its trace measure \(\mu=T\wedge\omega\) satisfies
\[
 \lim_{r\downarrow0}\sup_{x\in X}\frac{\mu(B(x,r))}{r^2}=0 .
\]
Here balls can be taken in the fixed smooth Hermitian metric.
\end{lemma}
\begin{proof}
We include the intrinsic proof of \cite[Proposition~2.1]{Oct}, based
on \cite[Lemma~5.3]{CX}, since its uniform density estimate is needed
below. The separation argument belongs to the theory of Levi currents
of Sibony \cite{Sibony} and of minimal kernels introduced by
S{\l}odkowski and Tomassini \cite{ST}. For compact local maximum sets
and Levi currents see Bianchi and Mongodi
\cite[Theorem~1.1 and Lemma~2.8]{BM}. Let \(\Sigma\) be the
locus where every continuous psh exhaustion fails to be strictly psh on a
neighborhood. By \cite[Lemma 2.1]{MS}, some half-line \((t_0,\infty)\)
lies in \(\rho(\Sigma)\). For each \(s>t_0\) there is no smooth strictly
psh function on a neighborhood of the entire compact level \(K_s=\{\rho=s\}\).
Indeed, normalize such a function \(v\) to take values in \((0,1)\) on
a compact band \(|\rho-s|\leq\delta\). Put \(g(t)=t+\tfrac12|t-s|\).
Use \(\max(g(\rho),\rho+\epsilon v)\) on the band and \(g(\rho)\)
elsewhere, with \(0<\epsilon<\delta/4\). The resulting continuous psh
exhaustion is strictly psh near \(K_s\), contradicting \(K_s\cap\Sigma\ne
\varnothing\).

In the space of continuous real \((1,1)\)-forms on \(K_s\), separate
the open positive-definite cone from
\(\{(\D f)|_{K_s}: f\in C_c^\infty(X,\R)\}\) by Hahn--Banach.
The separating functional annihilates the subspace and is positive on
nonnegative forms. Its value on \(\omega\) is strictly positive, since
all forms on the compact level are bounded above and below by multiples
of \(\omega\). Normalizing gives a positive current \(T_s\) of mass one,
supported on \(K_s\), with \(\D T_s=0\).

Fix \(t_0<a<b\). For each \(n\), average the \(T_s\) at the \(n\)
midpoints of a uniform partition of \([a,b]\), and pass to a weak
subsequence. The limit \(T\) has mass one, is supported in the compact
band \(\{a\leq\rho\leq b\}\), and satisfies \(\D T=0\). On a relatively
compact coordinate ball write \(\beta=\D|z|^2\). For centers in a
smaller ball, Skoda's local Lelong--Jensen monotonicity at a fixed
larger radius \cite[Proposition~1]{Skoda}, in the form recorded in
\cite[Appendix B, (B.1)--(B.2)]{DNS}, gives
\(\int_{B(x,r)}T_s\wedge\beta\leq C r^2\) uniformly in \(x,s\).
Only those midpoints in the range of \(\rho\) on \(B(x,r)\) can
contribute. Their proportion is at most
\(\osc_{B(x,r)}\rho/(b-a)+2/n\). Passing to the weak limit on the open
ball gives
\[
  \int_{B(x,r)}T\wedge\beta
  \leq \frac{Cr^2}{b-a}\osc_{B(x,r)}\rho .
\]
Uniform continuity on a compact neighborhood of the band, finitely many
smaller coordinate balls, and comparison of metrics give the assertion.
\end{proof}

\section{Smoothing a current with a Hermitian metric}

The density and Gaussian estimates below come from
\cite[Proposition~4.1]{Oct}, which uses K\"ahler heat flow.
Ni and Tam \cite{NiTam} use the Hodge heat equation to construct
Poincar\'e--Lelong potentials on complete K\"ahler manifolds.
For regularization of positive $dd^c$-closed currents by differences
of smooth positive $dd^c$-closed forms on a compact K\"ahler manifold,
see Dinh and Sibony \cite[Theorem~4.1]{DS}. Here support in a
prescribed neighborhood and a uniform lower bound for the approximating
forms must hold simultaneously.

\begin{lemma}\label{lem:heat}
Let \(T\) be a positive \((1,1)\)-current on a complex surface, with
\(\D T=0\), compact support \(K\), and zero Lelong numbers everywhere.
Fix a smooth Hermitian form \(\omega\) near \(K\). In any prescribed
relatively compact neighborhood of \(K\) there is a compact set
\(F\) supporting real smooth \((1,1)\)-forms \(A_t\), \(0<t<1\),
with
\[
 A_t\longrightarrow T,\qquad
 \sup_{0<t<1}\|A_t\|_{L^1}<\infty,\qquad
 \|\D A_t\|_{C^k}=O_{k,N}(t^N),\qquad
 \liminf_{t\downarrow0}\int A_t\wedge A_t\geq0.
\]
Here \(k,N\) are arbitrary nonnegative integers, the first
convergence is in currents, and the error estimates are as \(t\downarrow0\).
In fact \(A_t\geq-\eta_t\omega\) on \(F\) for some \(\eta_t\to0\).
\end{lemma}

\begin{proof}
Write \(L f=g^{i\bar j}\partial_i\partial_{\bar j}f\), the complex
trace of \(\D f\) for \(\omega\). We construct a real operator \(P\)
on Hermitian \((1,1)\)-forms satisfying \(P\D=\D L\).
In a holomorphic chart write \(U=iU_{k\bar l}\,dz^k\wedge d\bar z^l\).
Set
\begin{align*}
 (P_{\rm loc}U)_{k\bar l}
 ={}& L U_{k\bar l}
 +(\partial_k g^{i\bar j})\,\partial_{\bar j}U_{i\bar l}
 +(\partial_{\bar l}g^{i\bar j})\,\partial_iU_{k\bar j}\\
 &+(\partial_k\partial_{\bar l}g^{i\bar j})U_{i\bar j}.
\end{align*}
Substitution of \(U_{k\bar l}=\partial_k\partial_{\bar l}f\)
verifies the identity. Transport these local operators as operators
on the bundle and take a smooth partition-of-unity weighted sum.
This gives \(P\D=\D L\) on a neighborhood of \(K\), with scalar
principal symbol equal to that of \(L\).

The connection of this Laplace-type operator preserves the cone of
positive Hermitian forms under parallel transport. Indeed, in the
displayed chart the coefficient of each real first derivative in
\(P_{\rm loc}-L\) acts on a Hermitian matrix by
\(U\mapsto B U+ U B^*\). Indeed set
\((M_j)_k{}^i=\partial_k g^{i\bar j}\) and write \(z^j=x^j+iy^j\).
The coefficients of \(\partial_{x^j}\) and \(\partial_{y^j}\) are,
respectively,
\[
 U\longmapsto (M_jU+UM_j^*)/2,\qquad
 U\longmapsto \big((iM_j)U+U(iM_j)^*\big)/2 .
\]
Such endomorphisms form the Lie algebra of the congruence action.
An additional real scalar drift also has this form. Changing the
holomorphic frame acts by congruence, whose logarithmic derivative
belongs to the same Lie algebra. The construction of the connection
of a scalar-symbol Laplace-type operator multiplies these first-order
coefficients by the inverse principal symbol and a real scalar and
adds scalar multiples of the identity. Thus the connection matrices
remain in this Lie algebra, also after the weighted sum. Parallel
transport has the form \(U\mapsto H U H^*\), with \(H\) invertible.

For the heat construction, choose a relatively compact open set with
smooth boundary containing \(K\) in the given neighborhood. A finite
union of components is allowed. Double its closure as a
smooth manifold. Double a smooth vector bundle extending the
\((1,1)\)-bundle on this domain, using a collar trivialization in the
normal direction to the boundary. On the closed double extend the
principal metric, the bundle connection and the zeroth-order
endomorphism, keeping \(P,L\) and the real bundle on a fixed smaller
neighborhood of \(K\). This extension is only a smooth elliptic
construction. Let \(B_t=(e^{tP})'T\), the distributional transpose
on dual bundle densities. In the smaller neighborhood identify such
densities with \((1,1)\)-forms by the wedge pairing. In complex
dimension two the positive cone is its own dual for this pairing:
testing a Hermitian matrix against the rank-one positive forms tests
its value in all complementary complex directions.

We record the estimates used for this heat operator. On the closed
double, the heat parametrix of a Laplace-type operator, including an
operator with nonmetric connection, has leading kernel equal to
\((4\pi t)^{-2}\exp[-d(x,y)^2/(4t)]\) times parallel transport and a
smooth positive density factor. Here \(d\) is the metric distance
of its real principal symbol. The next terms are this Gaussian
times positive powers of \(t\) with smooth bounded bundle
coefficients on a fixed neighborhood of the diagonal. A smooth
cutoff with values in \([0,1]\) localizes this expansion. By
increasing its truncation order, the remainder is \(O(t^N)\) in
any prescribed \(C^k\) norm on the double. These are the
parametrix and remainder in \cite[Theorem~2.30]{BGV}. See also
\cite[Theorem~3.1]{Ludewig} for the near-diagonal
estimate for a general Laplace-type operator on a compact manifold.
In particular, \(B_t\) and its derivatives are \(O(t^N)\) on every
closed set separated from \(K\), for every \(N\).

Choose a nonnegative cutoff \(\chi\leq1\) equal to one near \(K\)
and supported in the unchanged neighborhood, and put
\(A_t=\chi B_t\), extended by zero. We check the error in
\(\D A_t\). On this neighborhood transposition of \(P\D=\D L\)
gives \(\D P'B=L'\D B\), where \(L'\) acts on scalar densities.
Consequently the compactly supported scalar density
\(r_t=\D(\chi B_t)\), viewed on the double, satisfies
\[
    (\partial_t-L')r_t=e_t,\qquad r_{0}=0.
\]
Indeed \(e_t=\D[\chi,P']B_t\), with the commutator
\([\chi,P']=\chi P'-P'\chi\). All its terms are supported where derivatives of
\(\chi\) are nonzero and contain only finitely many derivatives of
\(B_t\). The off-diagonal estimate gives
\(\|e_t\|_{C^k}=O_{k,N}(t^N)\); the initial datum is zero in
distributions because \(\D T=0\). Duhamel's formula implies the
same estimate for \(r_t\). To justify the zero initial term, first
start the formula at \(s>0\). For fixed \(t>0\),
\(e^{(t-s)L'}r_s\to0\) distributionally as \(s\downarrow0\),
by weak convergence of \(r_s\) and smooth convergence of the heat
operator on test functions at positive time. On the closed double the
scalar heat semigroup is bounded over \(0\leq t\leq1\) on each
fixed Sobolev space; apply the formula in sufficiently high Sobolev
order and then use Sobolev embedding. The heat approximation to
the identity also gives \(A_t\to T\).

Let \(\mu=T\wedge\omega\). Its total mass controls the variation of
all coefficients of \(T\). For some fixed \(C>0\) set
\[
 q_t(x)=t^{-2}\int_K \exp[-d(x,y)^2/(Ct)]\,d\mu(y).
\]
In the parametrix choose the diagonal neighborhood small enough
that geodesics with one endpoint in \(K\) lie in the unchanged
neighborhood. The transposed leading coefficient preserves positivity
by cone duality. The remaining finite list of coefficients is
bounded in norm by \(Ct\) times the Gaussian there. Thus
\[
       A_t=H_t+E_t,\qquad H_t\geq0,\qquad
       |H_t|\leq Cq_t,\qquad |E_t|\leq Ctq_t+O(t^2)
\]
on the fixed compact support, where \(H_t\) denotes the transposed
leading term times \(\chi\). In particular \(\|A_t\|_1\) is
uniformly bounded since \(\|q_t\|_1\leq C\mu(K)\).

For clarity, the Lelong assumption gives the uniform estimate
\[
   \sup_x\mu(B(x,r))=o(r^2).
\]
Here one may restrict centers to a fixed compact neighborhood of \(K\).
In each holomorphic coordinate ball the Euclidean trace mass
divided by \(r^2\) is nondecreasing with \(r\), by the local
Lelong--Jensen formula for positive \(\D\)-closed currents.
If the uniform assertion failed, there would be centers \(x_j\to x\)
and radii \(r_j\downarrow0\) with normalized Euclidean mass
bounded below. For every fixed sufficiently small \(s>0\),
monotonicity at \(x_j\) gives mass \(\geq c s^2\) in \(B(x_j,s)\);
this ball lies in \(B(x,2s)\) for large \(j\).
This contradicts \(\nu(T,x)=0\). A finite atlas and local uniform
comparability give the assertion in the chosen metric.
For a direct estimate set \(m(r)=\sup_x\mu(B(x,r))/r^2\).
This function tends to zero at zero and is bounded for all \(r>0\)
by compactness of the support and finiteness of the mass. Splitting
the Gaussian into the ball of radius \(\sqrt t\) and dyadic annuli gives
\[
 t\|q_t\|_\infty\le m(\sqrt t)+
 \sum_{k\ge0}4^{k+1}e^{-4^k/C}m(2^{k+1}\sqrt t)\longrightarrow0
\]
by dominated convergence. In particular
\(\|E_t\|_\infty=o(1)\), and the positivity of the transported
leading term yields \(A_t\geq-\eta_t\omega\) on \(F\) with
\(\eta_t\to0\). Integrating the pointwise inequality
\((A_t+\eta_t\omega)^2\geq0\) over \(F\) gives
\[
   \int_X A_t^2\ \geq\
   -2\eta_t\int_F A_t\wedge\omega-\eta_t^2\int_F\omega^2
   \ \geq\ -C\eta_t-C\eta_t^2 .
\]
Only the leading transport was required to preserve positivity.
\end{proof}

\section{A Dirichlet potential on continuous sublevels}

\begin{lemma}\label{lem:almost}
On a neighborhood of any fixed compact set, a continuous psh function
\(\rho\) has smooth approximations \(r_\epsilon\) converging uniformly
and satisfying \(\D r_\epsilon\geq-\epsilon\omega\).
\end{lemma}
\begin{proof}
Apply B{\l}ocki and Ko{\l}odziej \cite[Theorem~2]{BK} with
\(\gamma=0\) on a relatively compact open neighborhood of the given
set. The continuous \(\rho\) is locally bounded and has zero Lelong
numbers. Their smooth approximants decrease to \(\rho\) with Hessians
bounded below by \(-\epsilon_j\omega\), where \(\epsilon_j\downarrow0\).
Dini's theorem gives uniform convergence on a smaller compact
neighborhood. For the continuous case the same approximation with a
strictly negative Levi lower bound follows from Greene and Wu
\cite[\S4, Theorem~4.1 and Corollary~2]{GW}.
\end{proof}

\begin{lemma}[A Dirichlet potential on an exhaustion sublevel]\label{lem:dirichlet}
Let $(X,\omega)$ be a connected noncompact Hermitian surface with a continuous psh
exhaustion $\rho$. Let $U$ be a connected component of $\{\rho<c\}$,
and let $K\Subset U$ be compact. Suppose $T\ne0$ is a positive
$(1,1)$-current supported in $K$. Assume there are smooth real
$(1,1)$-forms $A_t$ supported in a fixed compact subset $K'\Subset U$,
converging to $T$ as currents as $t\downarrow0$, such that
\[
 \sup_{0<t<t_0}\|A_t\|_{L^1}<\infty,\qquad
 \|\D A_t\|_{C^0}\longrightarrow0,\qquad
 \liminf_{t\downarrow0}\int_X A_t\wedge A_t\ge0 .
\]
There is a psh function $F:U\to[-\infty,0)$ with $\D F=T$.
It is finite and smooth pluriharmonic on $U\setminus K$.
It tends uniformly to zero at $\partial U$.
\end{lemma}

\begin{proof}
The October~29 note \cite[Theorem~9.1 and Corollary~9.5]{Oct}
uses a global Green potential assuming completeness, infinite volume,
bounded K\"ahler curvature and decay of the positive Green kernel at
infinity. The discussion in \cite[\S9]{Oct}, drawing on Napier and
Ramachandran \cite[proof of Theorem~2.5]{NR}, already points to
hyperbolic sublevels and asks how their local separators can be made
compatible. We solve on relatively compact domains. The extra term to
estimate is the boundary integral of the square of the complex Hessian.
We use the boundary identity in Weber's K\"ahler argument
\cite[Lemmas~2.1--2.2]{Weber} for the complex trace of a Hermitian metric.

We use $d^c=\tfrac{i}{2}(\bar\partial-\partial)$, so
$dd^c=\D$. Set $L=\Lambda_\omega\D$ and $dV=\omega^2/2$.
This real scalar operator is elliptic with smooth coefficients and
zero zeroth-order term. Choose a relatively compact, connected,
smoothly bounded domain $W$ containing $\{\rho\le c+1\}$.
Increasing $W$ slightly if needed, fix a compact neighborhood $H$
of $\{\rho\le c+1/2\}$ in $W$. All the level surfaces used below
lie in the interior of $H$.

By Lemma~\ref{lem:almost}, on a neighborhood of $\overline W$ there
are smooth functions $s_j$ such that $s_j\to\rho$ uniformly on
$\overline W$ and $\D s_j\ge-\delta_j\omega$, where
$\delta_j\downarrow0$. Solve $Lp=1$ on $W$, with zero boundary
data; the classical Dirichlet theorem gives
$p\in C^\infty(\overline W)$. For $r_j=s_j+2\delta_jp$ we have
$Lr_j\ge0$, because $\Lambda_\omega\omega=2$. There is a sequence
$\epsilon_j>0$ tending to zero for which
\[
 \|r_j-\rho\|_{C^0(\overline W)}\le\epsilon_j,\qquad
 \D r_j\ge-\epsilon_j\omega\quad\hbox{on }\overline W.
\]
On $H$ we have a uniform energy bound
$\int_H |dr_j|^2\,dV\le C$. Indeed, take a nonnegative cutoff
$\eta\in C_c^\infty(W)$ equal to one on $H$, and a fixed $B$ making
$r_j+B>0$ on $W$. The product rule and ellipticity give
$L(r_j+B)^2\ge C_0|dr_j|^2$ on $\supp\eta$, with $C_0>0$.
The integral of $\eta L(r_j+B)^2$ is uniformly bounded by integration
against $L^*\eta$.

Write $h_j=\sqrt{\epsilon_j}$. By the coarea formula and Sard's
theorem, there is a regular value $c_j\in(c+h_j,c+2h_j)$ such that
\begin{equation}\label{eq:dirichlet-flux}
 \int_{\{r_j=c_j\}\cap W}|dr_j|\,d\sigma\le C/h_j .
\end{equation}
Here $d\sigma$ and the gradient are for the Riemannian metric of
$\omega$. Indeed the coarea integral over this interval is bounded
by the energy on $H$ for all large $j$. Since $\rho>c+1$ on
$\partial W$, these levels are compactly contained in $W$.
For all large $j$, $U\subset\{r_j<c_j\}\cap W$. Denote by $V_j$
the component of this smooth sublevel containing $U$. Its closure
is contained in the fixed compact set $H$.

Choose any sequence $t_j\downarrow0$ and abbreviate $A_j=A_{t_j}$.
On $V_j$ solve the smooth classical Dirichlet problem
\[
 Lu_j=\Lambda_\omega A_j,\qquad u_j|_{\partial V_j}=0.
\]
The solutions satisfy
\begin{equation}\label{eq:dirichlet-Green-bounds}
 \int_{V_j}|u_j|\,dV\le C,\qquad
 \sup_{V_j\cap H'}|u_j|\le C(H')
 \quad\text{for each compact }H'\subset H\setminus K'.
\end{equation}
For clarity, the constants are independent of the varying domains.
Use the positive Dirichlet Green kernel $G_W(x,y)$ of $-L$ on the
fixed smooth $W$, relative to $dV(y)$. The maximum principle gives
\[
 |u_j(x)|\le\int_{K'} G_W(x,y)|\Lambda_\omega A_j(y)|\,dV(y).
\]
The kernel is smooth off the diagonal, hence bounded on disjoint
compact sets. To bound its integral in the first variable, take
$0\leq\eta\in C_c^\infty(W)$ equal to one on $H$. The transpose of
the Dirichlet inverse of $-L$ applied to $\eta$ is smooth in the
interior by elliptic regularity. Its value at $y$ bounds
$\int_H G_W(x,y)\,dV(x)$ from above, uniformly for $y\in K'$.
This argument also allows the first-order drift of the Hermitian
operator. The uniform $L^1$ bound on $A_j$ proves
\eqref{eq:dirichlet-Green-bounds}.

Fix $d>0$ so small that $\rho<c-3d$ on $K\cup K'$.
Both $u_j$ and $-u_j$ are $L$-harmonic on the collar
$V_j\cap\{c-d<r_j<c_j\}$. On its inner boundary their absolute
values are bounded independently of $j$ by
\eqref{eq:dirichlet-Green-bounds}. The function $C(r_j-c_j)$ is
$L$-subharmonic. Comparison on the collar, taking $C$ large enough,
therefore yields
\begin{equation}\label{eq:dirichlet-collar-bound}
 |u_j|\le C(c_j-r_j) \quad\text{on }V_j\cap\{r_j>c-d\}.
\end{equation}
The inner level need not be regular for this maximum-principle
argument. At $\partial V_j$ write $du_j=b_j\,dr_j$.
The one-sided derivative of \eqref{eq:dirichlet-collar-bound} in an
inward transverse direction gives $|b_j|\le C$.

Following the trace correction of Buchdahl
\cite[\S2, especially Lemma~4]{Buchdahl} and its complete K\"ahler
counterpart \cite[Theorem~3.3 and Lemma~4.2]{Oct}, the smooth residual
$B_j=A_j-\D u_j$ is primitive on $V_j$, so
$B_j^2=-|B_j|_\omega^2dV$. Stokes' theorem gives
\[
 \int_{V_j}(\D u_j)^2
   =\int_{\partial V_j} d^cu_j\wedge\D u_j
   =\int_{\partial V_j}b_j^2 d^c r_j\wedge\D r_j.
\]
For the second equality, pull $d^cu_j=b_jd^c r_j$ back to
$\partial V_j$ and differentiate there. With the boundary oriented
outwards, $d^c r_j\wedge\omega$ restricts to a nonnegative volume
form bounded by $C|dr_j|\,d\sigma$. Only the complex tangential
Levi form enters $d^c r_j\wedge\D r_j$, whence
\[
 \int_{V_j}(\D u_j)^2
   \ge -C\epsilon_j \int_{\partial V_j}|dr_j|\,d\sigma
   \ge -C\sqrt{\epsilon_j}.
\]
Since $A_j$ is supported in $K'\Subset V_j$, integration by parts
has no boundary contribution in the cross term:
\[
 \bigg|\int_{V_j} A_j\wedge\D u_j\bigg|
   =\bigg|\int_{V_j} u_j\D A_j\bigg|
   \le C\|\D A_j\|_{C^0}\longrightarrow0 .
\]
The sign of each estimate now enters through the identity
\[
 0\leq\int_{V_j}|A_j-\D u_j|_\omega^2\,dV
 =-\int_X A_j^2+2\int_{V_j}u_j\D A_j-\int_{V_j}(\D u_j)^2 .
\]
The right side has nonpositive upper limit, proving strong convergence
of the residual to zero in $L^2(V_j)$.

By \eqref{eq:dirichlet-Green-bounds} a subsequence of $u_jdV$
converges weakly as measures on $U$. Regard its limit as a scalar
distribution $F$ via the fixed volume $dV$. The $L^2$ estimate and
the convergence $A_j\to T$ give $\D F=T$ on $U$. The distribution
has a psh representative. On $U\setminus K$ it is smooth by
elliptic regularity for $L$, and is pluriharmonic there.
Test both inequalities in \eqref{eq:dirichlet-collar-bound} against
nonnegative smooth compactly supported functions on
$U\cap\{\rho>c-d/2\}$. Uniform convergence of $r_j$ and weak convergence
of $u_j\,dV$ give $|F|\le C(c-\rho)$ almost everywhere there.
Both sides are continuous on this collar, giving the pointwise estimate.
Every boundary point of the component $U$ has $\rho=c$, so the
convergence to zero is uniform at $\partial U$. The psh maximum
principle gives $F\le0$ on $U$. An interior value zero would make
$F$ constant on the connected set $U$, contrary to $\D F=T\ne0$.
\end{proof}

\section{Exhaustion by Levi-flat boundaries}

\begin{lemma}[Gluing pluriharmonic collars on a surface]\label{lem:annuli}
Let $X$ be a connected noncompact complex surface. Suppose that every
compact subset of $X$ is contained in a relatively compact connected
domain with smooth boundary. Suppose each boundary component has a
single pluriharmonic defining function on a full two-sided neighborhood
of that component, with nonvanishing differential there.
Then $X$ admits a smooth plurisubharmonic exhaustion.
\end{lemma}

\begin{proof}
First fill the relatively compact complementary components of each such
domain $W$. More precisely, adjoin to $\overline W$ all relatively
compact components of $X\setminus\overline W$ and take the interior.
The compact smooth hypersurface $\partial W$ has finitely many connected
components. Its exterior collar along any one of them is connected, and
every component of $X\setminus\overline W$ meets such a collar. Thus
only finitely many components are filled. The result is a relatively
compact connected smooth domain, its boundary consists of components of
$\partial W$, and every component of the complement of its closure is
not relatively compact in $X$. Recursively applying the hypothesis to
compact sets containing the previous filled closure and a cofinal
sequence of compact sets gives an exhaustion by filled domains $U_n$,
$n\geq 1$, with $\overline U_n\subset U_{n+1}$.

Write $S_n=\partial U_n$. On a collar of $S_n$ take a defining
function $r_n$ as in the hypothesis, choosing its sign on each component
so that it increases outwards. Choose disjoint collars and sufficiently
small $\epsilon_n>0$. Moving $S_n$ to the levels
$S_n^-=\{r_n=-\epsilon_n\}$ and $S_n^+=\{r_n=\epsilon_n\}$ within
its collar gives domains $U_n^-$ and $U_n^+$, respectively, with
\[
 \overline U_n^-\subset U_n\subset\overline U_n\subset U_n^+,
 \qquad \overline U_n^+\subset U_{n+1}^-.
\]
These small collar moves are ambient isotopies supported in relatively
compact sets. Hence the perturbed domains are connected and
$X\setminus\overline{U_n^\pm}$ has no
relatively compact components. The perturbed boundaries are Levi-flat
because they are regular levels of a pluriharmonic function.

On $A_n=U_{n+1}^+\setminus\overline U_n^-$ solve the scalar Dirichlet
problem
\[
 \operatorname{tr}_{g}(\D h_n)=0,\qquad
 h_n|_{S_n^-}=0,\qquad h_n|_{S_{n+1}^+}=1,
\]
using any smooth Hermitian metric $g$ and $\D=i\partial\bar\partial$.
Each component of $A_n$ meets both boundary sets in its closure. It
meets the inner boundary by connectivity of $U_{n+1}^+$, as one sees by
joining a point of the component to $U_n^-$ inside $U_{n+1}^+$.
If its closure missed the outer boundary, it would be a relatively
compact component of $X\setminus\overline U_n^-$, which was excluded
above. Classical elliptic Dirichlet theory and the strong maximum
principle give a solution smooth up to the boundary and $0<h_n<1$ in
$A_n$.

We use Weber's harmonic separator argument
\cite[Propositions~1.1 and 3.2]{WeberEarly},
\cite[Lemma~1.1 and its proof]{Weber}. In the K\"ahler case one uses
ordinary harmonic functions. For a general Hermitian metric the
Dirichlet operator chosen above ensures directly that $\D h_n$
has trace zero. Its two eigenvalues are therefore
opposites, so $(\D h_n)^2$ is nonpositive and vanishes only where
$\D h_n=0$. Use $d^c=\frac{i}{2}(\bar\partial-\partial)$, so that
$dd^c=\D$. On each boundary component $h_n$ is constant, and its Levi
form on the complex tangent line is its normal derivative times the
Levi form of a local defining function, up to the normalization of the
normal. Thus it vanishes on that line. Since $d^ch_n$ annihilates the
same real two-plane, $d^ch_n\wedge\D h_n$ restricts to zero on
$\partial A_n$. Stokes' theorem gives
\[
 \int_{A_n}(\D h_n)^2=\int_{\partial A_n}d^ch_n\wedge\D h_n=0.
\]
It follows that $h_n$ is pluriharmonic on $A_n$.

The overlapping collars permit gluing using only comparisons of values.
Set $F_0=0$ on $U_1^+$ and $b_1=-1$. Recursively put
\[
 F_n=b_n+a_n h_n,\qquad
 m_n=\min_{S_{n+1}^-}h_n>0,\qquad
 b_{n+1}=b_n+\tfrac12 a_n m_n.
\]
Choose $a_n>0$ large enough that $a_nm_n/2\geq1$ and
$F_n>F_{n-1}$ on $S_n^+$. The second requirement is possible since
$\min_{S_n^+}h_n>0$ and $F_{n-1}$ is already defined up to $S_n^+$.
On $S_{n+1}^-$ the opposite comparison at the next transition follows
from
\[
 F_n\geq b_n+a_nm_n>b_{n+1}=F_{n+1}.
\]
At the initial inner boundary $F_0=0>b_1=F_1$. On the overlap
$U_n^+\setminus\overline U_n^-$ use a smooth convex regularized
maximum of $F_{n-1}$ and $F_n$, nondecreasing in both variables.
Compactness and the two strict boundary comparisons allow its smoothing
parameter to be chosen so that it equals $F_{n-1}$ near $S_n^-$ and
$F_n$ near $S_n^+$. It can also be chosen to dominate the ordinary
maximum. This is the regularized-maximum gluing used in psh smoothing
(see \cite{Richberg} and \cite[\S1]{BK}). Use the corresponding single function off the overlaps. The
result is smooth and plurisubharmonic. Since $F_n\geq b_n$ and
$b_{n+1}\geq b_n+1$, it tends to infinity off compact sets.
\end{proof}

\begin{proof}[Proof of Theorem~\ref{thm:main}]
A constant suffices when \(X\) is compact. If \(X\) is a modification of
a Stein space, pull back a smooth strictly psh exhaustion of that space
by the proper modification. In all other cases fix \(T\) from
Lemma~\ref{lem:shell}.

To see cofinality explicitly, place any prescribed compact set and a
neighborhood of \(\supp T\) inside a compact connected set \(H\). Such a
set can be obtained by covering the compact sets with finitely many
relatively compact coordinate balls and joining them by paths. Choose
\(c>\max_H\rho\) and take the component of \(\{\rho<c\}\) containing
\(H\). Inside this component choose a relatively compact open
neighborhood of \(\supp T\) and apply Lemma~\ref{lem:heat} there.
The forms it gives have common compact support in the component, as
required by Lemma~\ref{lem:dirichlet}. That lemma gives a strictly negative psh function
\(u\), tending uniformly to zero at the boundary and pluriharmonic
off \(\supp T\). Its upper semicontinuity gives \(\sup_H u<0\).
Choose a regular value \(s\) between \(\sup_H u\) and zero for the
smooth function off \(\supp T\), by Sard's theorem. All points of
\(\{u=s\}\) are then in the smooth locus, and this level is compact
by the boundary estimate. The component of \(\{u<s\}\) containing
\(H\) has smooth Levi-flat boundary and contains the prescribed compact
set. The same \(u-s\) defines the boundary on a full two-sided
neighborhood where it is pluriharmonic and its differential is nonzero.
Lemma~\ref{lem:annuli} completes
the construction on \(X\).
\end{proof}

\section{Normal complex surfaces}

\begin{corollary}\label{cor:normal}
Let \(Y\) be a connected normal complex surface with a continuous psh
exhaustion. Then \(Y\) admits a smooth psh exhaustion, with smoothness
understood in local ambient embeddings.
\end{corollary}
\begin{proof}
The compact case is immediate. Take a proper resolution
\(\pi:X\to Y\) which is an isomorphism over the regular locus and has
simple normal crossings exceptional divisor \cite{BierMil}. Normal surface singularities
are isolated, the fibers of \(\pi\) are connected, and
\(\pi_*\mathcal O_X=\mathcal O_Y\) by normality (holomorphic functions
upstairs are locally bounded downstairs by properness and extend across
the singular set). Connectedness of the fibers also follows from Stein
factorization. Use the pulled-back exhaustion
\(\rho=\rho_Y\circ\pi\) on \(X\).

If \(X\) is a modification of a Stein space \(Z\), the proper map
\(q:X\to Z\) is constant on the compact connected fibers of \(\pi\).
It induces a continuous map \(\bar q:Y\to Z\), since \(\pi\) is a
proper surjection. In local embeddings of \(Z\), descent of the
coordinate functions by \(\pi_*\mathcal O_X=\mathcal O_Y\) makes
\(\bar q\) holomorphic. It is proper, since
\(\bar q^{-1}(J)=\pi(q^{-1}(J))\) for compact \(J\).
Pull back a smooth strictly psh exhaustion of \(Z\).

In the other case make one additional choice in the proof of
Theorem~\ref{thm:main}. A component \(V\) of \(\{\rho<c\}\) contains
every exceptional fiber it meets, by the choice of \(\rho\) and
connectedness of the fibers. The psh potential of
Lemma~\ref{lem:dirichlet} is constant on each such fiber, possibly
with value \(-\infty\). This follows from the maximum principle on
the normalization of each irreducible compact curve and connectedness.
There are only finitely many singular points beneath \(\overline V\),
by properness of the exhaustion and discreteness of the closed singular
locus. In choosing the regular level \(s\) above \(\sup_H u\) we may
therefore avoid the values on all these fibers. The smooth boundary
constructed in the theorem is disjoint from the entire exceptional
locus. Filling relatively compact complementary components retains
this property. In Lemma~\ref{lem:annuli}, shrink every collar to have
closure disjoint from the exceptional locus. The family of collars
is locally finite by the nested exhaustion. Consequently the final
smooth psh exhaustion \(v\) of \(X\) is constant or pluriharmonic
on an open neighborhood of each connected exceptional fiber.

For descent we use the normalization of local holomorphic primitives
on a compact curve from Mongodi, S{\l}odkowski and Tomassini
\cite[proof of Lemma~4.1]{MST}. We give the argument for a reducible
exceptional fiber, allowing critical points of the pluriharmonic function.
Let \(E=\pi^{-1}(p)\)
be an exceptional fiber, and \(a=v|_E\) its constant value. A sufficiently
small coordinate neighborhood of any point of \(E\) meets \(E\) in a
connected set consisting of one smooth branch or two crossing branches.
There \(v-a=\operatorname{Re} f\) for a holomorphic function \(f\).
On every branch the real part is zero, so \(f\) is constant there.
Adjust its imaginary constant so that \(f|_E=0\) in that neighborhood.
These local choices agree as germs at all common points of \(E\):
their difference is a purely imaginary constant locally and vanishes
at that point. They glue to a holomorphic function on a neighborhood
of \(E\). Indeed, take a finite cover by smaller charts with closures
inside the original ones. Overlap points where two functions disagree
cannot accumulate on \(E\) in the closures of the smaller charts, by
the germ agreement just proved. Shrink around \(E\) to remove all
such overlaps. Properness of \(\pi\) then supplies a neighborhood
\(B\) of \(p\) whose full inverse image lies in the gluing neighborhood.
The holomorphic function descends on \(B\), so \(v\) descends there as
the real part of a holomorphic function plus \(a\). This is smooth
and psh in the ambient sense. Away from the singular points descent
uses the isomorphism. The descended function is exhaustive because
its closed sublevels are the images of compact sublevels of \(v\).
\end{proof}

For arbitrary reduced complex surfaces, the cuspidal construction in
\cite[Proposition~5.7]{CX}
already gives a negative answer. Its divisor complement carries a
continuous psh exhaustion, and every ambient \(C^1\) psh function on
it is constant.

\begin{thebibliography}{MST18}

\bibitem[BGV92]{BGV} N. Berline, E. Getzler and M. Vergne,
\emph{Heat Kernels and Dirac Operators}, Springer, 1992,
\href{https://doi.org/10.1007/978-3-642-58088-8}{Chapter 2}.

\bibitem[BM22]{BM} F. Bianchi and S. Mongodi,
On minimal kernels and Levi currents on weakly complete complex manifolds,
\emph{Proc. Amer. Math. Soc.} \textbf{150} (2022), 3927--3939.
\href{https://arxiv.org/abs/2102.05328}{Author version}.

\bibitem[BM97]{BierMil} E. Bierstone and P. D. Milman,
Canonical desingularization in characteristic zero by blowing up the
maximum strata of a local invariant, \emph{Invent. Math.}
\textbf{128} (1997), 207--302.
\href{https://arxiv.org/abs/alg-geom/9508005}{Author version}.

\bibitem[B{\l}K07]{BK} Z. B{\l}ocki and S. Ko{\l}odziej,
On regularization of plurisubharmonic functions on manifolds,
\emph{Proc. Amer. Math. Soc.} \textbf{135} (2007), 2089--2093.
\href{https://gamma.im.uj.edu.pl/~blocki/publ/appr.pdf}{Theorem~2}.

\bibitem[Buc99]{Buchdahl} N. Buchdahl, On compact K\"ahler surfaces,
\emph{Ann. Inst. Fourier} \textbf{49} (1999), 287--302.
\href{https://www.numdam.org/article/AIF_1999__49_1_287_0.pdf}{\S2, Lemma~4}.

\bibitem[CX26]{CX} Z. Chen and B. Xiao,
\emph{Regularity of plurisubharmonic exhaustions and related problems},
manuscript, version~8, November 20, 2026.
Theorem numbers refer to that version.
\href{https://github.com/chen-zhijian/psh-exhaustion}{Project repository}.

\bibitem[DNS22]{DNS} T.-C. Dinh, V.-A. Nguy\^en and N. Sibony,
Unique ergodicity for foliations on compact K\"ahler surfaces,
\emph{Duke Math. J.} \textbf{171} (2022), 2627--2698.
\href{https://arxiv.org/abs/1811.07450}{Author version, Appendix B}.

\bibitem[DS04]{DS} T.-C. Dinh and N. Sibony,
Regularization of currents and entropy, \emph{Ann. Sci. \'Ecole
Norm. Sup.} (4) \textbf{37} (2004), 959--971.
\href{https://arxiv.org/abs/math/0407532}{Author version, Theorem~4.1}.

\bibitem[GW79]{GW} R. E. Greene and H. Wu,
$C^\infty$ approximations of convex, subharmonic, and plurisubharmonic
functions, \emph{Ann. Sci. \'Ecole Norm. Sup.} (4) \textbf{12} (1979),
47--84. \href{https://doi.org/10.24033/asens.1361}{\S4}.

\bibitem[Lud19]{Ludewig} M. Ludewig, Strong short-time asymptotics and
convolution approximation of the heat kernel,
\emph{Ann. Global Anal. Geom.} \textbf{55} (2019), 371--394.
\href{https://arxiv.org/abs/1607.05152}{Author version, Theorems 1.1 and 3.1}.

\bibitem[MS20]{MS} S. Mongodi and Z. S{\l}odkowski,
Domains with a continuous exhaustion in weakly complete surfaces,
\emph{Math. Z.} \textbf{296} (2020), 1011--1019.
\href{https://arxiv.org/abs/1904.03497}{Author version, Lemma 2.1}.

\bibitem[MST18]{MST} S. Mongodi, Z. S{\l}odkowski and G. Tomassini,
Weakly complete complex surfaces,
\emph{Indiana Univ. Math. J.} \textbf{67} (2018), 899--935.
\href{https://arxiv.org/abs/1504.07215}{Author version, Lemma~4.1}.

\bibitem[Nak80]{Nakano} S. Nakano, problem in I. Wakabayashi (ed.),
\emph{Collected problems around several complex variables} (in Japanese),
\emph{Sugaku} \textbf{32} (1980), 161--187.
\href{https://doi.org/10.11429/sugaku1947.32.161}{\S I.1, Problem~1, p.~162}.

\bibitem[NR97]{NR} T. Napier and M. Ramachandran, The Bochner--Hartogs
dichotomy for weakly 1-complete K\"ahler manifolds, \emph{Ann. Inst.
Fourier} \textbf{47} (1997), 1345--1365.
\href{https://doi.org/10.5802/aif.1602}{Theorem~2.5}.

\bibitem[NT13]{NiTam} L. Ni and L.-F. Tam, Poincar\'e--Lelong equation via
the Hodge--Laplace heat equation, \emph{Compos. Math.} \textbf{149}
(2013), 1856--1870. \href{https://doi.org/10.1112/S0010437X12000322}{Publisher}.

\bibitem[Ohs81]{Ohsawa} T. Ohsawa, Weakly 1-complete manifold and Levi problem,
\emph{Publ. Res. Inst. Math. Sci.} \textbf{17} (1981), 153--164.
\href{https://doi.org/10.2977/PRIMS/1195186709}{Proposition~1.4}.

\bibitem[RN26]{Oct} \emph{Intrinsic current methods for the K\"ahler surface
regularization problem}, unpublished manuscript, October~29, 2026.
The byline reads \emph{Research note}.

\bibitem[Ric68]{Richberg} R. Richberg, Stetige streng pseudokonvexe Funktionen,
\emph{Math. Ann.} \textbf{175} (1968), 257--286.
\href{https://doi.org/10.1007/BF02063212}{Publisher}.

\bibitem[Sib18]{Sibony} N. Sibony, Levi problem in complex manifolds,
\emph{Math. Ann.} \textbf{371} (2018), 1047--1067.
\href{https://doi.org/10.1007/s00208-017-1539-x}{Publisher}.

\bibitem[Sko82]{Skoda} H. Skoda, Prolongement des courants positifs, ferm\'es
de masse finie, \emph{Invent. Math.} \textbf{66} (1982), 361--376.
\href{https://doi.org/10.1007/BF01389217}{Publisher}.

\bibitem[ST04]{ST} Z. S{\l}odkowski and G. Tomassini, Minimal kernels of
weakly complete spaces, \emph{J. Funct. Anal.} \textbf{210} (2004),
125--147. \href{https://doi.org/10.1016/S0022-1236(03)00182-4}{Publisher}.

\bibitem[TV08]{TV} G. Tomassini and V. V\^aj\^aitu, Weakly 1-complete
surfaces with singularities and applications, \emph{Michigan Math. J.}
\textbf{57} (2008).
\href{https://doi.org/10.1307/mmj/1231770355}{Publisher}.

\bibitem[V{\^\i}j22]{Vij} V. V\^{\i}j\^{\i}itu, Holomorphic convexity and
pseudoconvex surfaces, \emph{Ann. Sc. Norm. Super. Pisa Cl. Sci.} (5)
\textbf{23} (2022), 1659--1678.
\href{https://doi.org/10.2422/2036-2145.202011_018}{Theorem~1.1}.

\bibitem[Web19]{WeberEarly} B. Weber, Topology of K\"ahler manifolds with
weakly pseudoconvex boundary, \emph{Michigan Math. J.} \textbf{68}
(2019), 727--742. \href{https://arxiv.org/abs/1110.4571}{Author version}.

\bibitem[Web22]{Weber} B. Weber, Pseudoconvex submanifolds in K\"ahler
4-manifolds, \emph{Commun. Contemp. Math.}
\href{https://doi.org/10.1142/S0219199724500147}{2450014}.
\href{https://arxiv.org/abs/2208.00975}{Preprint (2022), Lemmas~1.1, 2.1 and 2.2}.

\end{thebibliography}
\end{document}
